\subsection{Runtime and scalability}\label{sec:runtime}
\owner{Wei Yew, Lam}

\todo{Wei Yew}{D8, all 234{,}294 sales: K-means under 0.06\,s; HDBSCAN 212\,s; DBSCAN not run
(neighbour lists need about 6\,GB). Plot runtime against $n$ (\nolinkurl{weiyew/sc4020-hdb/experiments/scalability.py}).}
% Lam: lam/results/fashion_mnist_results.csv (fit_seconds; single fit at the chosen parameters, tuning excluded)
D5 was sampled down from 70{,}000 to 3{,}000 images because three methods scale badly with $n$ in
50 dimensions. Spectral clustering needs an $n\times n$ affinity matrix: dense, at $n=70{,}000$ it
takes about 39\,GB in double precision, and the eigenvector step grows faster than linearly even with
a sparse graph. DBSCAN and HDBSCAN rely on neighbour queries, and in 50 dimensions spatial indexes
prune little, so these queries approach the $O(n^2)$ brute-force cost. On the 3{,}000-point sample
one fit takes 0.011\,s for DBSCAN, 0.079\,s for K-Means++, 0.18\,s for Spectral, 0.30\,s for HDBSCAN
and 2.80\,s for GMM. GMM is the slowest here: with full covariances each EM step costs
$O(nkd^2)$ with $d=50$. These times exclude parameter search, which refits DBSCAN 15 times and
HDBSCAN 6 times.
\todo{Ky}{Existing Fig.~9 runtime at $n\leq23{,}412$, if kept.}

\begin{figure}[h]\centering
\figtodo{Wei Yew}{runtime vs number of sales, log--log, one line per method}
\caption{\todo{Wei Yew}{caption}}
\label{fig:runtime}
\end{figure}
